Este capítulo é reservado à análise dos dados do instrumento descrito na Seção~\ref{sec:protocolo-avaliacao}. Até o fechamento desta revisão a coleta não tinha sido feita: o formulário está escrito e revisado, mas não foi distribuído, e o encaminhamento ao comitê de ética da instituição continua pendente. Por isso não há resultados de participantes aqui.

As seções abaixo definem desde já o que será relatado e com quais critérios. Se essa definição fosse feita depois de ver os dados, seria possível escolher o recorte mais favorável à solução; defini-la agora evita essa escolha. Os critérios de cálculo dos escores estão na Seção~\ref{sec:protocolo-avaliacao} e não são repetidos aqui.

% =====================================================================
% A PREENCHER APÓS A COLETA.
% Cada seção diz o que relatar e com qual critério. Manter a ordem e os
% critérios definidos aqui; se algo mudar depois de ver os dados, a
% mudança precisa ser declarada no texto.
% =====================================================================

\section{Composição do Painel}
\label{sec:resultado-painel}

% A PREENCHER: quantos responderam; distribuição por titulação, área de
% atuação e tempo de experiência; quantos lecionam ou já lecionaram; quais
% ferramentas de ensino o painel já conhecia.
% Numa avaliação feita por especialistas, o perfil de quem avalia faz parte
% do resultado: apresentar antes dos escores.
% Registrar o que os respondentes disseram ter feito na plataforma e o que
% personalizaram: sem isso não dá para saber se uma avaliação negativa veio
% de quem usou a ferramenta ou de quem não experimentou o recurso avaliado.

\section{Avaliação do Núcleo de Tradução}
\label{sec:resultado-nucleo}

% A PREENCHER: escores do bloco de qualidade técnica, item a item, com
% atenção ao item que diz se o código intermediário exibido corresponde ao
% que o programa executa.
% Relatar separadamente quantos marcaram "não sei" em cada item: num bloco
% dirigido a quem tem formação na área, muitas respostas assim indicam que
% o recurso não foi usado, e não que foi mal avaliado.
% Reunir aqui também os itens sobre a exibição das fases intermediárias
% (tabela de tokens, código intermediário e execução passo a passo) que
% vieram do bloco de avaliação dos recursos.

\section{Suficiência do Subconjunto da Linguagem}
\label{sec:resultado-subconjunto}

% A PREENCHER: resultado do item sobre o subconjunto ser suficiente para
% uma disciplina introdutória, cruzado com as construções marcadas como
% faltantes. Se o painel convergir para a mesma construção ausente, isso
% delimita até onde vai o que foi implementado e ajuda a priorizar os
% trabalhos futuros no núcleo de tradução.

\section{Mensagens de Erro com Vocabulário Personalizado}
\label{sec:resultado-diagnosticos}

% A PREENCHER: item sobre clareza das mensagens de erro no bloco de
% qualidade técnica, o item correspondente no bloco de avaliação dos
% recursos e o item pedagógico sobre mostrar erros no vocabulário do
% próprio aluno.
% Comparar com os relatos da questão aberta sobre mensagens confusas: cada
% relato é um caso concreto em que a mensagem não se adaptou ao vocabulário
% em uso, e deve ser rastreado até a fase que a emitiu.

\section{Usabilidade e Experiência}
\label{sec:resultado-usabilidade}

% A PREENCHER: escore SUS geral, com intervalo de confiança só se houver
% vinte respondentes ou mais; com menos, apresentar como indicativo e não
% afirmar que há diferença entre subgrupos. Comparar com a média de
% referência e com as faixas de interpretação.
% UEQ-S: relatar as dimensões pragmática e hedônica separadas, na faixa de
% -3 a +3, comparando com os valores de referência do instrumento.

\section{Potencial Pedagógico e Intenção de Uso}
\label{sec:resultado-pedagogico}

% A PREENCHER: bloco de potencial pedagógico (oito itens em terceira
% pessoa) e bloco vindo do TAM.
% Discutir se a opinião do painel concorda com a premissa de que concentrar
% a flexibilidade da linguagem na análise léxica traz ganho pedagógico.

\section{Riscos e Limitações Apontados pelo Painel}
\label{sec:resultado-riscos}

% A PREENCHER: bloco de riscos, relatado SEPARADAMENTE e avisando que ali
% concordar significa apontar um problema, e não aprovar.
% Não somar nem tirar média junto com o bloco de potencial pedagógico: as
% escalas apontam em sentidos opostos.
% Comparar cada risco com as decisões de projeto descritas neste capítulo,
% dizendo quais já estão resolvidas pela implementação atual e quais
% continuam em aberto.

\section{Respostas Abertas}
\label{sec:resultado-abertas}

% A PREENCHER: análise das cinco questões abertas.
% Citar trechos apenas de quem autorizou a citação, sem identificação,
% conforme o item de autorização do formulário.

\section{Síntese e Confronto com os Objetivos}
\label{sec:resultado-sintese}

% A PREENCHER: retomar os objetivos específicos e dizer, para cada um, o
% que a avaliação sustenta e o que continua sem verificação.
% Manter a distinção da Subseção \ref{subsec:limites-interpretacao}: o que
% foi coletado é a opinião de profissionais sobre a plausibilidade da
% redução da barreira, e não a prova dessa redução.
